% Adapted from academic-drawing/templates/tikz_template.tex to the book styles.
% Main ports: south -> north. The dashed loop carries a new subquestion.
\begin{tikzpicture}[
  every node/.style={font=\figurefont\small,align=center},
  block/.style={book node,text width=64mm,minimum height=12mm},
  note/.style={text width=51mm,align=left,text=BookMuted},
  flow/.style={book arrow,rounded corners=1.2mm},
  retry/.style={flow,dashed,draw=BookAmber}
]
  \node[block,fill=FigureInput!12,draw=FigureInput] (question)
    at (-22mm,0) {问题与访问协议\\资料版本、时点、预算};
  \node[block,fill=FigureInput!12,draw=FigureInput] (candidates)
    at (-22mm,-23mm) {候选资料\\召回A、B，排除干扰C};
  \node[block,fill=FigureModel!12,draw=FigureModel] (evidence)
    at (-22mm,-46mm) {选定证据\\保留来源、对象与限定条件};
  \node[block,fill=FigureModel!12,draw=FigureModel] (answer)
    at (-22mm,-69mm) {答案陈述与来源标记\\时间引用A，截止条件引用B};
  \node[block,fill=FigureOutput!12,draw=FigureOutput] (check)
    at (-22mm,-92mm) {引用核验与交付\\支持则保留，不足则限定或弃答};

  \draw[flow] (question.south) -- (candidates.north);
  \draw[flow] (candidates.south) -- (evidence.north);
  \draw[flow] (evidence.south) -- (answer.north);
  \draw[flow] (answer.south) -- (check.north);
  \node[note] at (48mm,-23mm) {检索遗漏\\只召回B，缺活动时间};
  \node[note] at (48mm,-46mm) {组合错误\\把报名截止当成活动时间};
  \node[note] at (48mm,-69mm) {无支持引用\\引用A解释校外报名条件};
  \node[note] at (48mm,-92mm) {独立验收\\答案、证据覆盖、引用支持};
  \draw[retry] (evidence.west) -- (-66mm,-46mm)
    -- (-66mm,0) -- (question.west);
  \node[font=\figurefont\footnotesize,text=BookAmber,rotate=90]
    at (-72mm,-23mm) {未决子问题：预算内再查询};
\end{tikzpicture}
